\section{Details on Existence of RARE}
\label{sec:exist-extra}

\subsection{Security and Overpricing Parameters}
\label{sec:exist-extra:market-params}

Let $\evec_j$ denote the $j\Th$ standard basis vector in $\mathbb{R}^m$.
We first isolate the utility estimate used in the security and overpricing arguments.

\begin{lemma}[Uniform utility improvement]
\label{thm:uniform-utility-improvement}
Let $v: [0, 1] \times [0, 1]^m \to \mathbb{R}$ be continuous, and suppose that
$v(\eps, \cdot)$ is strictly monotone for every $\eps \in [0, 1]$.
Fix a good $k$ and $\eta \in (0, 1)$.
There exists $\delta > 0$ such that, for every $\eps \in [0, 1]$ and
$\zvec, \yvec \in [0, 1]^m$,
\[ z_k \le 1 - \eta \quad \text{and} \quad
    \norm{\yvec - (\zvec + \eta\evec_k)}_\infty < \delta
    \implies v(\eps, \yvec) > v(\eps, \zvec). \]
\end{lemma}
\begin{proof}
By continuity, strict monotonicity, and compactness,
\[ \gamma \defeq \min_{\substack{\eps \in [0, 1], \; \zvec \in [0, 1]^m \\ z_k \le 1 - \eta}}
    \bigl(v(\eps, \zvec + \eta\evec_k) - v(\eps, \zvec)\bigr) > 0. \]
Uniform continuity of $v$ gives $\delta > 0$ such that changing the bundle
by less than $\delta$ in the infinity norm changes its value by less than $\gamma/2$,
uniformly over $\eps$.
Thus, under the stated conditions,
$v(\eps, \yvec) > v(\eps, \zvec + \eta\evec_k) - \gamma/2 > v(\eps, \zvec)$.
\end{proof}

\thmSecurity*
\begin{proof}
Fix seller $j$ and use buyer 1 to secure positive revenue.
Apply \cref{thm:uniform-utility-improvement} with
$v(\eps, \zvec) \defeq u_1(\zvec)$, $k = j$, and $\eta = 1/4$, obtaining $\delta_j > 0$.
Set
\[ \sigma_j \defeq b_1 \min(1, \delta_j), \qquad
   \rho_j \defeq \frac{\sigma_j}{4}, \qquad
   \lambda_j \defeq \frac{\sigma_j}{4b_1}. \]
In particular, $0 < \lambda_j \le 1/4$ and $\lambda_j < \delta_j$.

Fix prices with $q_j = \sigma_j$ and an allocation $\xvec \in \Dem(\qvec)$.
Suppose that $\sum_i x_{i, j} < 1/4$.
Replace only buyer 1's bundle by
\[ \yvec_1 \defeq (1 - \lambda_j)\xvec_1 + \tfrac{1}{4}\evec_j. \]
This bundle lies in $[0, 1]^m$, and its cost is at most
\[ \qvec^T\yvec_1 \le (1 - \lambda_j)b_1 + \frac{\sigma_j}{4} = b_1. \]
In the rivalrous setting, the modification also preserves the stock constraints:
the total allocation of good $j$ is less than $1/2$, and every other total weakly decreases.
Since $x_{1, j} < 1/4$ and
\[ \norm{\yvec_1 - (\xvec_1 + \tfrac{1}{4}\evec_j)}_\infty
    = \lambda_j\norm{\xvec_1}_\infty < \delta_j, \]
\cref{thm:uniform-utility-improvement} gives $u_1(\yvec_1) > u_1(\xvec_1)$.
This contradicts buyer 1's optimality in the non-rivalrous setting
and Pareto optimality in the rivalrous setting.
Hence $\sum_i x_{i, j} \ge 1/4$, so $R_j(\qvec, \xvec) \ge \sigma_j/4 = \rho_j$.
\end{proof}

\thmOverpricing*
\begin{proof}
Fix seller $j$, and let
\[ a_j \defeq \frac{\rho_j}{2\kappa}, \qquad
   \eta_j \defeq \min\biggl(\frac{1}{2}, \frac{a_j}{2nB}\biggr). \]
For each buyer $i$ and good $k \ne j$, apply \cref{thm:uniform-utility-improvement}
to $v(\eps, \zvec) \defeq u_i(\zvec)$ with increment $\eta_j$ in good $k$.
Let $\delta_j > 0$ be the minimum of the resulting finitely many constants, and set
\[ L_j \defeq \max\biggl(B, \sigma_j, \frac{B}{\delta_j}\biggr). \]

Fix prices and demand $\xvec$ as in the statement, and suppose that
$R_j(\qvec, \xvec) > a_j$.
Choose a good $k \ne j$ with $q_k = B$, and a buyer $i$ with
$q_jx_{i, j} > a_j/n$.
We will replace buyer $i$'s purchases of good $j$ by $\eta_j$ units of good $k$.

First, there is enough room for this increase in good $k$.
In the non-rivalrous setting, buyer $i$'s budget constraint gives
\[ Bx_{i, k} \le b_i - q_jx_{i, j} < B - a_j/n,
    \qquad \text{so} \qquad x_{i, k} < 1 - 2\eta_j. \]
In the rivalrous setting, summing the budget constraints gives
\[ B \sum_{\ell = 1}^n x_{\ell, k}
    \le B - R_j(\qvec, \xvec) < B - a_j. \]
Thus the unsold quantity of good $k$ is greater than $a_j/B \ge 2n\eta_j$.
In either setting, adding $\eta_j$ units of good $k$ to buyer $i$ is feasible
with respect to the consumption and stock constraints.

Now replace only buyer $i$'s bundle by
\[ \yvec_i \defeq \xvec_i - x_{i, j}\evec_j + \eta_j\evec_k. \]
Its cost decreases, since $q_jx_{i, j} > a_j/n$ and $B\eta_j \le a_j/(2n)$.
Moreover,
\[ \norm{\yvec_i - (\xvec_i + \eta_j\evec_k)}_\infty
    = x_{i, j} \le \frac{b_i}{q_j} \le \frac{B}{q_j} < \delta_j. \]
By \cref{thm:uniform-utility-improvement}, this strictly improves buyer $i$'s utility.
As in the security proof, this contradicts demand in either setting.
Consequently, $R_j(\qvec, \xvec) \le a_j$.
\end{proof}

\begin{restatable}[Robust parameters]{lemma}{thmRobustParams}
\label{thm:robust-params}
Let $\Ical \defeq ([n], [m], (u_i)_{i=1}^n, (b_i)_{i=1}^n)$ be a market instance.
For every $\eps \in [0, 1]$, define the perturbed instance
$\Ical^{(\eps)} \defeq ([n], [m], (u_i^{(\eps)})_{i=1}^n, (b_i)_{i=1}^n)$, where
\[ u_i^{(\eps)}(\xvec) \defeq u_i(\xvec) + \eps \sum_{k=1}^m x_k(2-x_k). \]
%
There exist constants $\rho_j > 0$, $\sigma_j \ge \rho_j/\nu$, and $L_j \ge \max(B, \sigma_j)$
for every seller $j \in [m]$, independent of $\eps$,
such that for every $\eps \in [0, 1]$, seller $j \in [m]$, price vector $\qvec \in \mathbb{R}_{\ge 0}^m$,
and allocation $\xvec \in \Dem(\qvec; \Ical^{(\eps)})$, both of the following hold:
\begin{tightenum}
\item (security) $q_j = \sigma_j \implies R_j(\qvec, \xvec) \ge \rho_j$.
\item (overpricing)
    $\displaystyle q_j > L_j \text{ and } \max_{k \in [m] \setminus \{j\}} q_k = B
        \implies R_j(\qvec, \xvec) \le \frac{\rho_j}{2\kappa}$.
\end{tightenum}
\end{restatable}
\begin{proof}
For every buyer $i$, the function
\[ v_i(\eps, \zvec) \defeq u_i(\zvec) + \eps \sum_{k = 1}^m z_k(2 - z_k) \]
is jointly continuous in $(\eps, \zvec)$ and strictly monotone in $\zvec$.
Indeed, $z \mapsto z(2 - z)$ is increasing on $[0, 1]$, and $u_i$ is strictly monotone.
Thus \cref{thm:uniform-utility-improvement} applies uniformly over $\eps \in [0, 1]$.

For each seller $j$, choose $\delta_j$ using that lemma with $v_1$, good $j$,
and increment $1/4$, and set $\sigma_j = b_1 \min(1, \delta_j)$ and $\rho_j = \sigma_j/4$.
The feasible bundle modification in the proof of \cref{thm:security}
then gives the security bound for every $\eps \in [0, 1]$ with these same constants.

Next, set $a_j = \rho_j/(2\kappa)$ and
$\eta_j = \min(1/2, a_j/(2nB))$.
Apply \cref{thm:uniform-utility-improvement} to every $v_i$ and good $k \ne j$
with increment $\eta_j$, and let $\widehat{\delta}_j > 0$ be the minimum of its constants.
Set $L_j = \max(B, \sigma_j, B/\widehat{\delta}_j)$.
The feasible bundle modification in the proof of \cref{thm:overpricing}
now gives the overpricing bound for every $\eps \in [0, 1]$.
All of these parameters are independent of $\eps$.
\end{proof}

\subsection{RARE for Continuous Allocation Rule}
\label{sec:exist-extra:continuous}

\begin{lemma}
\label{thm:superlevel-closed}
Let $C \subseteq \mathbb{R}^m$ be a closed set and let $f: C \to \mathbb{R}$ be continuous.
Then the set $S = \{\xvec \in C: f(\xvec) \ge 0\}$ is closed.
\end{lemma}
\begin{proof}
Since $\mathbb{R}^m$ is a metric space, it suffices to show that
$S$ contains the limits of all its convergent sequences.

Let $(\xvec^{(k)})_{k \in \N}$ be a sequence in $S$ with $\xvec^{(k)} \to \xvec$
for some $\xvec \in \mathbb{R}^m$. We must show $\xvec \in S$.

First, each $\xvec^{(k)} \in S \subseteq C$, and $C$ is closed, so $\xvec \in C$.
In particular, $f(\xvec)$ is defined.

Second, since $f$ is continuous on $C$ and $(\xvec^{(k)})$ is a sequence in $C$
converging to $\xvec \in C$, we have $f(\xvec^{(k)}) \to f(\xvec)$.
Each $f(\xvec^{(k)}) \ge 0$, and $[0, \infty)$ is closed in $\mathbb{R}$,
so $f(\xvec) = \lim_{k \to \infty} f(\xvec^{(k)}) \ge 0$.

Hence $\xvec \in C$ and $f(\xvec) \ge 0$, i.e., $\xvec \in S$. Therefore $S$ is closed.
\end{proof}

\subsection{Limits of Affordable Allocations and Maximizers}
\label{sec:exist-extra:limits}

The following two lemmas will be used for both non-rivalrous and rivalrous demand.

\begin{lemma}[Limits of affordable allocations]
\label{thm:affordable-limits}
For each price vector $\pvec \in \mathbb{R}_{\ge 0}^m$ and buyer $i$, let
$K_i(\pvec) \defeq \{\zvec \in [0, 1]^m: \pvec^T\zvec \le b_i\}$.
Let $\mathcal{F}(\pvec)$ be either $\prod_{i=1}^n K_i(\pvec)$ or $\AffR(\pvec)$.
These sets are nonempty, compact, and convex.
If $\pvec^{(t)} \to \pvec$, then:
\begin{tightenum}
\item If $\xvec^{(t)} \in \mathcal{F}(\pvec^{(t)})$ and $\xvec^{(t)} \to \xvec$,
    then $\xvec \in \mathcal{F}(\pvec)$.
\item For every $\xvec \in \mathcal{F}(\pvec)$, there are
    $\xvec^{(t)} \in \mathcal{F}(\pvec^{(t)})$ with $\xvec^{(t)} \to \xvec$.
\end{tightenum}
\end{lemma}
\begin{proof}
Both sets contain the zero allocation and are bounded sets defined by closed linear inequalities,
so they are compact and convex.
The first claim follows by taking limits in these inequalities.

For the second claim, fix $\xvec \in \mathcal{F}(\pvec)$ and set
\[ \alpha_i^{(t)} \defeq
    \frac{b_i}{\max\{b_i, (\pvec^{(t)})^T\xvec_i\}},
    \qquad \xvec_i^{(t)} \defeq \alpha_i^{(t)}\xvec_i. \]
Each $\alpha_i^{(t)}$ lies in $(0, 1]$, and $\xvec_i^{(t)}$ is affordable at $\pvec^{(t)}$.
Scaling bundles down preserves the consumption bounds and, in the rivalrous setting,
the stock constraints. Thus $\xvec^{(t)} \in \mathcal{F}(\pvec^{(t)})$.
Since $\pvec^T\xvec_i \le b_i$ and $b_i > 0$, we have $\alpha_i^{(t)} \to 1$,
and hence $\xvec^{(t)} \to \xvec$.
\end{proof}

\begin{lemma}[Limits of maximizers]
\label{thm:maximizer-limits}
Let $\mathcal{F}$ be either of the feasible sets in \cref{thm:affordable-limits}.
Let $V_t, V: [0, 1]^{n \times m} \to \mathbb{R}$ be continuous functions such that
$V_t \to V$ uniformly.
If $\pvec^{(t)} \to \pvec$, $\xvec^{(t)} \to \xvec$, and
\[ \xvec^{(t)} \in \argmax_{\yvec \in \mathcal{F}(\pvec^{(t)})} V_t(\yvec), \]
then $\xvec \in \argmax_{\yvec \in \mathcal{F}(\pvec)} V(\yvec)$.
\end{lemma}
\begin{proof}
By \cref{thm:affordable-limits}, $\xvec \in \mathcal{F}(\pvec)$.
Fix any $\yvec \in \mathcal{F}(\pvec)$, and use the same lemma to choose
$\yvec^{(t)} \in \mathcal{F}(\pvec^{(t)})$ with $\yvec^{(t)} \to \yvec$.
Optimality gives $V_t(\xvec^{(t)}) \ge V_t(\yvec^{(t)})$.
Uniform convergence and continuity of $V$ imply that the two sides converge to
$V(\xvec)$ and $V(\yvec)$, respectively.
Thus $V(\xvec) \ge V(\yvec)$, as required.
\end{proof}

\subsection{RARE for Non-Rivalrous Goods}
\label{sec:exist-extra:non-rival}

\thmUniqueNrDemand*
\begin{proof}
Let $V(\xvec) \defeq \sum_{i=1}^n u_i(\xvec_i)$.
Maximizing $V$ over $\prod_i K_i(\pvec)$ is equivalent to maximizing each buyer's utility
over her affordable bundles, so its set of maximizers is exactly $\DemNR(\pvec)$.
The objective is continuous and strictly concave, and the feasible set is nonempty,
compact, and convex by \cref{thm:affordable-limits}.
Consequently, there is a unique maximizer $A(\pvec)$.

To prove continuity, let $\pvec^{(t)} \to \pvec$.
The allocations $A(\pvec^{(t)})$ lie in the compact set $[0, 1]^{n \times m}$.
By \cref{thm:maximizer-limits} with $V_t = V$, every convergent subsequence of these
allocations has limit in $\DemNR(\pvec) = \{A(\pvec)\}$.
It follows that the whole sequence converges to $A(\pvec)$, proving continuity.
\end{proof}

\begin{lemma}[Demand survives limits]
\label{thm:demand-limit}
Let $\Ical \defeq ([n], [m], (u_i)_{i=1}^n, (b_i)_{i=1}^n)$ be a market instance.
For every $\eps \in [0, 1]$, define the perturbed instance
$\Ical^{(\eps)} \defeq ([n], [m], (u_i^{(\eps)})_{i=1}^n, (b_i)_{i=1}^n)$, where
\[ u_i^{(\eps)}(\xvec) \defeq u_i(\xvec) + \eps \sum_{j=1}^m x_j(2 - x_j). \]
If $(\eps_t)_{t \in \N} \to 0$, $(\qvec^{(t)})_{t \in \N} \to \qvec$, $(\xvec^{(t)})_{t \in \N} \to \xvec$,
and $\xvec^{(t)} \in \DemNR(\qvec^{(t)}; \Ical^{(\eps_t)})$, then $\xvec \in \DemNR(\qvec; \Ical)$.
\end{lemma}
\begin{proof}
Define
\[ V_t(\yvec) \defeq \sum_{i=1}^n u_i^{(\eps_t)}(\yvec_i),
    \qquad V(\yvec) \defeq \sum_{i=1}^n u_i(\yvec_i). \]
Since $0 \le z(2 - z) \le 1$ on $[0, 1]$, we have
$\sup_{\yvec \in [0, 1]^{n \times m}} |V_t(\yvec) - V(\yvec)| \le nm\eps_t \to 0$.
The budgets are unchanged by perturbation, and $\xvec^{(t)}$ maximizes $V_t$
over $\prod_i K_i(\qvec^{(t)})$.
By \cref{thm:maximizer-limits}, $\xvec$ maximizes $V$ over $\prod_i K_i(\qvec)$,
which is equivalent to $\xvec \in \DemNR(\qvec; \Ical)$.
\end{proof}

\subsection{RARE for Rivalrous Goods}
\label{sec:exist-extra:rival}

\begin{lemma}[Demand survives limits]
\label{thm:demand-limit-rival}
Let $\Ical \defeq ([n], [m], (u_i)_{i=1}^n, (b_i)_{i=1}^n)$ be a market instance.
Let $W: \mathbb{R}_{\ge 0}^n \to \mathbb{R}$ be a continuous, strictly monotone function.
For every $\eps \in [0, 1]$, define the perturbed instance
$\Ical^{(\eps)} \defeq ([n], [m], (u_i^{(\eps)})_{i=1}^n, (b_i)_{i=1}^n)$, where
\[ u_i^{(\eps)}(\xvec) \defeq u_i(\xvec) + \eps \sum_{j=1}^m x_j(2 - x_j). \]
If $(\eps_t)_{t \in \N} \to 0$, $(\qvec^{(t)})_{t \in \N} \to \qvec$, $(\xvec^{(t)})_{t \in \N} \to \xvec$,
and $\xvec^{(t)} \in \Dem^W(\qvec^{(t)}; \Ical^{(\eps_t)})$, then $\xvec \in \Dem^W(\qvec; \Ical)$.
\end{lemma}
\begin{proof}
Define
\begin{align*}
V_t(\yvec) & \defeq W\bigl(u_1^{(\eps_t)}(\yvec_1), \ldots, u_n^{(\eps_t)}(\yvec_n)\bigr), \\
V(\yvec) & \defeq W\bigl(u_1(\yvec_1), \ldots, u_n(\yvec_n)\bigr).
\end{align*}
For each $i$, $u_i^{(\eps_t)}$ converges uniformly to $u_i$, with error at most $m\eps_t$.
All utility vectors lie in the compact box $\prod_i [0, u_i(\vecOne) + m]$.
Uniform continuity of $W$ on this box therefore gives uniform convergence $V_t \to V$.
The feasible set $\AffR$ depends only on prices and budgets, so it is unchanged by perturbation.
Apply \cref{thm:maximizer-limits} to conclude that $\xvec$ maximizes $V$ over $\AffR(\qvec)$,
i.e., $\xvec \in \Dem^W(\qvec; \Ical)$.
\end{proof}

\thmUniqueRivalDemand*
\begin{proof}
Let $V(\xvec) \defeq W(u_1(\xvec_1), \ldots, u_n(\xvec_n))$.
We first show that $V$ is strictly concave.
Fix distinct allocations $\xvec, \yvec$ and $\lambda \in (0, 1)$, and write
\[ \avec \defeq (u_i(\xvec_i))_{i=1}^n, \qquad
    \bvec \defeq (u_i(\yvec_i))_{i=1}^n, \qquad
    \mathbf{c} \defeq (u_i(\lambda\xvec_i + (1 - \lambda)\yvec_i))_{i=1}^n. \]
Concavity of the utilities gives $\mathbf{c} \ge \lambda\avec + (1 - \lambda)\bvec$ coordinatewise.
Since $\xvec_i \ne \yvec_i$ for some buyer $i$, strict concavity of $u_i$
makes at least one of these inequalities strict.
Strict monotonicity and concavity of $W$ now give
\begin{align*}
V(\lambda\xvec + (1 - \lambda)\yvec)
    & = W(\mathbf{c}) > W(\lambda\avec + (1 - \lambda)\bvec) \\
    & \ge \lambda V(\xvec) + (1 - \lambda)V(\yvec).
\end{align*}
Thus $V$ is continuous and strictly concave.
Since $\AffR(\pvec)$ is nonempty, compact, and convex,
its maximizer $A(\pvec)$ is unique.
Strict monotonicity of $W$ ensures that this maximizer is Pareto optimal,
so $A$ is an allocation rule.

Finally, let $\pvec^{(t)} \to \pvec$.
By compactness of $[0, 1]^{n \times m}$ and \cref{thm:maximizer-limits} with $V_t = V$,
every convergent subsequence of $A(\pvec^{(t)})$ has limit equal to the unique maximizer $A(\pvec)$.
Hence $A(\pvec^{(t)}) \to A(\pvec)$, proving continuity.
\end{proof}

\thmRivalRareExists*
\begin{proof}
For $\eps \in [0, 1]$, let $\Ical^{(\eps)}$ have utilities
\[ u_i^{(\eps)}(\zvec) \defeq u_i(\zvec) + \eps \sum_{k = 1}^m z_k(2 - z_k). \]
Choose common parameters $\sigma_j, \rho_j, L_j$ using \cref{thm:robust-params}.
For $\eps_t = 1/t$, the perturbed utilities are strictly concave.
By \cref{thm:unique-rival-demand}, their $W$-rivalrous demand defines a unique continuous
allocation rule $A^{(\eps_t)}$.
Apply \cref{thm:rare-exists-continuous} with $\nu = \kappa = 1$ and $\beta = B$ to obtain
prices $\pvec^{(t)} \in B\Delta_m$ such that $p_j^{(t)} \ge \rho_j$ for every $j$.
Let $\xvec^{(t)} \defeq A^{(\eps_t)}(\pvec^{(t)})$.
Then $R_j(\pvec^{(t)}, \xvec^{(t)}) = p_j^{(t)}$ by \cref{thm:accounting}.

By compactness, pass to a joint convergent subsequence
\[ (\pvec^{(t)}, \xvec^{(t)}) \to (\pvec, \xvec)
    \in B\Delta_m \times [0, 1]^{n \times m}. \]
Relabel this subsequence and its perturbations using the same notation;
we still have $\eps_t \to 0$.
Fix this subsequence before choosing a seller or deviation.
The lower bound gives $p_j \ge \rho_j > 0$, so $\pvec \in B\Delta_m^\circ$.
By \cref{thm:demand-limit-rival}, $\xvec \in \Dem^W(\pvec; \Ical)$.
Continuity of revenue in prices and allocations gives $R_j(\pvec, \xvec) = p_j$ for every $j$.

Fix a seller $j$ and a deviation price $q_j \ge 0$.
For each $t$, \cref{thm:rare-exists-continuous} gives a response
$\qvec_{-j}^{(t)} \in \prod_{k \ne j} [p_k^{(t)}, L_k]$.
With $\qvec^{(t)} = (q_j, \qvec_{-j}^{(t)})$ and
$\zvec^{(t)} \defeq A^{(\eps_t)}(\qvec^{(t)})$, its revenue guarantees are
\[ R_j(\qvec^{(t)}, \zvec^{(t)}) \le p_j^{(t)}, \qquad
    R_k(\qvec^{(t)}, \zvec^{(t)}) \ge p_k^{(t)} \quad \text{for every }k \ne j. \]
The pairs $(\qvec_{-j}^{(t)}, \zvec^{(t)})$ lie in the compact set
$\bigl(\prod_{k \ne j} [0, L_k]\bigr) \times [0, 1]^{n \times m}$.
Take a further joint convergent subsequence with limit $(\qvec_{-j}, \zvec)$.
This further subsequence may depend on $j$ and $q_j$, but the limits
$\pvec, \xvec$ remain fixed and $\eps_t \to 0$ still holds.
We have $\qvec_{-j} \in \prod_{k \ne j} [p_k, L_k]$ and
$\zvec \in \Dem^W(\qvec; \Ical)$ by \cref{thm:demand-limit-rival}.
Taking limits in the revenue inequalities gives
\[ R_j(\qvec, \zvec) \le p_j = R_j(\pvec, \xvec), \qquad
    R_k(\qvec, \zvec) \ge p_k \quad \text{for every }k \ne j, \]
as required.
\end{proof}
